\section{AIM Master Architecture}

\begin{figure}[h]
\centering
\begin{tikzpicture}[
  node distance=1.25cm,
  box/.style={
    draw,
    rectangle,
    rounded corners,
    minimum width=5.8cm,
    minimum height=0.85cm,
    align=center,
    font=\small
  },
  smallbox/.style={
    draw,
    rectangle,
    rounded corners,
    minimum width=2.7cm,
    minimum height=0.75cm,
    align=center,
    font=\small
  },
  arrow/.style={->, thick},
  feedback/.style={->, thick, dashed},
  note/.style={font=\small, align=center}
]

% Main vertical pipeline
\node[box] (crs) {CRS / Georeferencing \\ $\mathcal{C}$};

\node[box, below=of crs] (world) {World Space \\ measured / designed coordinates};

\node[box, below=of world] (w2t) {World $\rightarrow$ Track Projection \\ $\mathcal{W}: x \mapsto (s,d)$};

\node[box, below=of w2t] (track) {Intrinsic Track Space \\ arc-length $s$, offset $d$};

% Hybrid model layer
\node[box, below=of track] (model) {Hybrid Alignment Model \\ $\kappa(s)=\kappa_{\mathrm{param}}(s;\mathbf{T})+\delta\kappa(s)$};

\node[smallbox, below left=0.9cm and 0.15cm of model] (param) {Parametric Layer \\ design intent};
\node[smallbox, below right=0.9cm and 0.15cm of model] (corr) {Correction Layer \\ measured deviations};

\node[box, below=1.55cm of model] (state) {Railway State \\ pose3: $(\gamma,\theta,\kappa,\phi)$};

\node[box, below=of state] (real) {Realization Operator \\ $\mathcal{R}$};

\node[box, below=of real] (realized) {Realized Geometry \\ built / observable track};

\node[box, below=of realized] (output) {World Output \\ visualization, validation, export};

% Main arrows
\draw[arrow] (crs) -- (world);
\draw[arrow] (world) -- (w2t);
\draw[arrow] (w2t) -- (track);
\draw[arrow] (track) -- (model);
\draw[arrow] (model) -- (state);
\draw[arrow] (state) -- (real);
\draw[arrow] (real) -- (realized);
\draw[arrow] (realized) -- (output);

% Hybrid internal arrows
\draw[arrow] (param.north) -- (model.south west);
\draw[arrow] (corr.north) -- (model.south east);

% Feedback loops
\draw[feedback]
  (realized.east) -- ++(1.6,0) |- node[pos=0.25, right]{Optimization / SQP} (model.east);

\draw[feedback]
  (output.west) -- ++(-1.6,0) |- node[pos=0.25, left]{Measurement feedback} (w2t.west);

% Notes
\node[note, right=1.6cm of real] (filternote)
  {$\mathcal{R}$ acts as\\low-pass filter};

\node[note, right=1.6cm of w2t] (lodnote)
  {LOD-dependent\\projection};

\node[note, left=1.6cm of model] (blocknote)
  {block structure:\\$\kappa$ / $\phi$ / corrections};

\draw[feedback] (filternote) -- (real);
\draw[feedback] (lodnote) -- (w2t);
\draw[feedback] (blocknote) -- (model);

\end{tikzpicture}

\caption{
Level-2 AIM master architecture.
The framework combines CRS handling, nonlinear world-to-track projection,
hybrid curvature modelling, pose3 state reconstruction, realization, and
optimization feedback. The central modelling object is not the geometric
curve itself, but the curvature-driven and realization-aware operator
system producing it.
}
\end{figure}
